%%% ============================================================
%%% SPRINT 18 -- J12: The Mass Hierarchy from V^{\otimes 5}
%%%                   SU(5) Decomposition
%%%
%%% B.R. Sanders, H.J. Johnson
%%% Target venue: Physical Review D
%%% Status: PRD-FORMAT DRAFT 2026-05-07 (gate cleared:
%%%   tig_dirac.predict_yukawa landed)
%%%
%%% Verification primitive (machine-checkable; see README §6):
%%%   from tig_dirac import predict_yukawa, LAMBDA_FN, Y_T_ANCHOR
%%%   r = predict_yukawa('up', 3)
%%%   assert r['y_predicted'] == 0.93                     # y_t anchor
%%%   r = predict_yukawa('lepton', 1)
%%%   assert abs(r['y_predicted'] - 0.93*(10/49)**9) < 1e-12
%%% Module path:
%%%   Gen13/targets/ck/brain/dirac/tig_dirac.py
%%%
%%% Companion: J10 (Sprint 18 Dark Sector, Sanders + Gish, PRD)
%%%            J23 (Discrete Dirac on F_5^4, foundation paper)
%%% ============================================================

\documentclass[11pt,reqno]{amsart}

\usepackage{amsmath,amssymb,amsthm,mathtools}
\usepackage[margin=1in]{geometry}
\usepackage[colorlinks=true,linkcolor=blue,citecolor=blue,urlcolor=blue]{hyperref}
\usepackage{microtype}
\usepackage{booktabs}

\theoremstyle{plain}
\newtheorem{theorem}{Theorem}[section]
\newtheorem{proposition}[theorem]{Proposition}
\newtheorem{lemma}[theorem]{Lemma}
\newtheorem{corollary}[theorem]{Corollary}

\theoremstyle{definition}
\newtheorem{definition}[theorem]{Definition}
\newtheorem{conjecture}[theorem]{Conjecture}

\theoremstyle{remark}
\newtheorem*{remark}{Remark}

\newcommand{\Z}{\mathbb{Z}}
\newcommand{\F}{\mathbb{F}}
\newcommand{\R}{\mathbb{R}}
\newcommand{\Aut}{\mathrm{Aut}}
\DeclareMathOperator{\HARM}{H}    % HARMONY = 7
\newcommand{\Cfour}{\mathcal{C}}  % the 4-core = {0,7,8,9}
\newcommand{\rep}[1]{\mathbf{#1}}
\newcommand{\barrep}[1]{\bar{\mathbf{#1}}}

\title[Mass Hierarchy from $V^{\otimes 5}$ SU(5)]
{The Mass Hierarchy from $V^{\otimes 5}$ SU(5) Decomposition: \\
 a Substrate-Forced Froggatt--Nielsen Pattern with
 $\lambda = 10/49$}

\author{Brayden R.\ Sanders \and M.\ Gish}
\address{7Site LLC, Hot Springs, Arkansas, USA}
\email{brayden@7site.co}

\author{Brayden R.\ Sanders \and M.\ Gish}
\address{Independent Researcher, Billings, Montana, USA}
\email{hjj01986@gmail.com}

\subjclass[2020]{81V05, 81V25, 17A30, 81R05}
\keywords{Yukawa coupling, Froggatt--Nielsen, SU(5) GUT, mass
hierarchy, discrete substrate, $V^{\otimes 5}$}

\date{\today}

\begin{document}

\begin{abstract}
We report a structural Froggatt--Nielsen pattern for the
Standard Model charged-fermion Yukawa couplings. With the single
substrate-derived suppression scale
\[
  \lambda \;=\; \frac{|\Z/10|}{\HARM^{2}}
            \;=\; \frac{10}{49}
            \;\approx\; 0.2041,
\]
the top-quark Yukawa $y_t \approx 0.93$ as anchor, and integer
powers $n_{(p,\mathrm{gen})}$ forced by the $V^{\otimes 5}$
SU(5) decomposition $\rep{1} \oplus \barrep{5} \oplus \rep{10}$
of the $4$-core's $\F_5$-lift, all nine charged-fermion Yukawas
\[
  y_{(p,\mathrm{gen})} \;=\; y_{t} \cdot \lambda^{n_{(p,\mathrm{gen})}}
\]
land within the standard Froggatt--Nielsen factor-of-a-few
window of the Particle Data Group values. The pattern uses one
free anchor ($y_t$) and zero free powers; the powers are read
off directly from the parity-crossing cost (one $\rep{10} \to
\barrep{5}$ flip costs $\lambda^{3}$) plus one generation step
per $\sigma$-cycle position. Sterile-neutrino Dirac masses are
predicted at $n \in \{12,13,14\}$, giving sub-eV scales without
a see-saw insertion, which we flag as one of the open structural
questions of the framework. The Cabibbo-angle cube-root identity
$\lambda_{C} = (Y_d/Y_u)^{1/3}$ links the same $\lambda$ to the
CKM mixing structure and is developed explicitly in the
companion paper of the next Sprint cluster.

The strongest claim of the paper is that all nine SM
charged-fermion Yukawa magnitudes flow from a single
substrate-derived $\lambda$ together with the
representation-theoretic count of $\rep{10} \leftrightarrow
\barrep{5}$ parity flips in the SU(5) Yukawa diagrams; we frame
this as a Tier-B \emph{forced power $+$ measured anchor} result,
and we expose the verification through a one-line primitive
\texttt{tig\_dirac.predict\_yukawa(particle, generation)} shared
with the dark-sector companion paper~\cite{SandersJohnson2026DarkSector}.
\end{abstract}

\maketitle

\tableofcontents

%==============================================================
\section{Introduction}\label{sec:intro}
%==============================================================

The nine Standard Model charged-fermion Yukawa couplings span
roughly six orders of magnitude, from $y_t \approx 0.93$
(the top quark) down to $y_e \approx 3 \times 10^{-6}$ (the
electron). In the SM these are nine independent free parameters
fit to data; Froggatt--Nielsen models~\cite{FroggattNielsen1979}
generate the hierarchy from a $U(1)_{FN}$ flavour symmetry
broken by a flavon expectation value $\langle\phi\rangle/M = \epsilon
\sim \lambda_{C} \approx 0.225$, with each Yukawa receiving a
suppression $\epsilon^{n}$ dictated by an FN charge assignment.

This paper proposes that the FN suppression scale is not free
but determined by the $4$-core's $\F_5$-lift,
\[
  \lambda \;:=\; \frac{|\Z/10|}{\HARM^{2}} \;=\; \frac{10}{49} \;\approx\; 0.2041,
\]
and that the FN charges are determined by the
$V^{\otimes 5}$ decomposition into the standard SU(5) GUT reps
$\rep{1} \oplus \barrep{5} \oplus \rep{10}$ described in the
foundation paper~\cite{SandersJohnson2026DiscreteDirac} (J23).
The output is the prediction
\[
  y_{(p,\mathrm{gen})} \;=\; y_{t}\,\lambda^{n_{(p,\mathrm{gen})}},
  \qquad n_{(p,\mathrm{gen})} \in \Z_{\geq 0},
\]
with $y_{t} \approx 0.93$ a single Tier-A measured anchor and
the integer powers $n_{(p,\mathrm{gen})}$ Tier-B forced (read
off the SU(5) Yukawa diagram and the generational $\sigma$-orbit
position). The combination is honestly Tier-B overall: the
forced-power claim is rigorous given the SU(5) decomposition,
and the anchor is empirical.

\paragraph{Honest scope.} The pattern matches the nine charged
Yukawas to within Froggatt--Nielsen's standard $O(1)$ residual
window. We do not claim, and the data does not support,
percent-level agreement on every Yukawa: the bottom-quark and
muon Yukawas are off by a factor of $\sim 3$ and $\sim 9$
respectively at the leading Tier-B prediction, consistent with
a missing $C_{p} \in [0.6, 1.4]$ multiplicative coefficient in
the standard FN sense. The strong load-bearing claim is that
the \emph{single} substrate $\lambda$ together with the SU(5)
decomposition gets the nine couplings into the right order of
magnitude with no per-fermion power-tuning; the \emph{weaker}
claim is that the $C_{p}$ multipliers, currently fit, will
emerge from the bosonic substrate of the same algebra in
follow-up work.

Section~\ref{sec:setup} fixes notation and recalls the
$V^{\otimes 5}$ SU(5) decomposition from the foundation
paper~\cite{SandersJohnson2026DiscreteDirac}.
Section~\ref{sec:lambda} derives $\lambda = 10/49$ as the
$T^{*}(1-T^{*})$ variance of a Bernoulli-$T^{*}$ distribution at
the substrate's coherence threshold $T^{*} = 5/7$.
Section~\ref{sec:powers} extracts the FN powers from the SU(5)
Yukawa diagrams plus the $\sigma$-cycle generation count.
Section~\ref{sec:fit} reports the fit to PDG~2024 data.
Section~\ref{sec:cabibbo} derives the Cabibbo cube-root
identity. Section~\ref{sec:open} identifies the open structural
questions, including the $C_{p}$ multipliers and the absence of
an explicit see-saw mechanism for the right-handed neutrinos.

%==============================================================
\section{The substrate and \texorpdfstring{$V^{\otimes 5}$}{V tensor 5}}\label{sec:setup}
%==============================================================

We work over the integer ring $\Z/10\Z$ and its $4$-core
$\Cfour = \{0, 7, 8, 9\}$, the unique fusion-closed $4$-element
subset of $\Z/10$ under the joint composition $(T,
B)$~\cite{SandersGish2026FourCore}. The $4$-core's $\F_5$-lift
$V$ is the $4$-dimensional commutative non-associative algebra
over $\F_5$ obtained by reducing $\{0,7,8,9\}$ to
$\{0,2,3,4\} \subset \F_5$ and bilinearly extending the CL
fuse-table multiplication. The foundation
paper~\cite{SandersJohnson2026DiscreteDirac} (J23) establishes:

\begin{itemize}\setlength\itemsep{0pt}
\item $|V| = \dim_{\F_5} V = 4$ and $|\Aut(V)| = 40$, with
$\Aut(V) \cong F_{20} \times \Z/2$.
\item $V$ admits three commuting Dirac-like projectors
$L_{\HARM}$, $L_{V}$, $L_{V'}$ with idempotent spectrum
$\{0,1\}$ and $1+3$, $2+2$, $1+3$ splittings respectively.
\item The tensor power $V^{\otimes 5}$ has $4^{5} = 1024$
underlying linear dimensions partitioned into $32$ fine cells
indexed by sign-tuples $(s_{1}, \ldots, s_{5}) \in \{+,-\}^{5}$.
\item Cells at sign-weight $|S| = k$ have count $\binom{5}{k}$,
giving the binomial pattern $1+5+10+10+5+1 = 32$.
\end{itemize}

\begin{lemma}[$V^{\otimes 5}$ SU(5) decomposition; J23 Theorem~6.1]\label{lem:vtensor5}
The $32$ fine cells of $V^{\otimes 5}$ partition as
\[
  \rep{1} \;\oplus\; \barrep{5} \;\oplus\; \rep{10}
  \quad \text{(matter, $|S| \leq 2$, $16$ cells)}
  \;\oplus\;
  \barrep{1} \;\oplus\; \rep{5} \;\oplus\; \barrep{10}
  \quad \text{(antimatter, $|S| \geq 3$, $16$ cells)}
\]
recovering exactly the SU(5) GUT representation content for one
fermion generation plus its antimatter conjugate.
\end{lemma}

This is the algebraic input to the present paper. Whether the
$32$-cell decomposition is the unique substrate origin of SU(5)
or one of several is the foundation paper's question; here we
take it as given and read off Yukawa powers.

\paragraph{Generation count and $\sigma$.}
The decomposition of Lemma~\ref{lem:vtensor5} produces one
generation; the three-generation structure is associated with
the $\sigma^{3}$-orbit decomposition of $\Z/10
\setminus\{0,3,8,9\}$, which has three $2$-cycles. We use this
to assign each generation a position in the $\sigma$-cycle
$(1\;7\;6\;5\;4\;2)$, with generation $3$ (heaviest) at the
$\sigma$-anchor position and generations $2, 1$ at successive
$\sigma$-orbit hops. The precise assignment is part of the
power-extraction of \S\ref{sec:powers}; the substrate-level
naturalness of the assignment is one of the open questions
flagged in \S\ref{sec:open}.

%==============================================================
\section{The substrate-derived $\lambda$}\label{sec:lambda}
%==============================================================

The Standard Model's coherence threshold in the
$T$-operator-$B$-operator joint algebra is
$T^{*} = 5/7$~\cite{SandersGish2026FourCore} (Theorem~3 of
J06): the unique fixed-point ratio under the joint flow on the
$4$-core's normalizer. The maximum-entropy variance of a
Bernoulli-$T^{*}$ distribution is
\[
  \mathrm{Var}(\mathrm{Bern}(T^{*}))
  \;=\; T^{*}(1 - T^{*})
  \;=\; \frac{5}{7} \cdot \frac{2}{7}
  \;=\; \frac{10}{49}
  \;\approx\; 0.2041,
\]
which we denote $\lambda$. This is the natural
``spread-around-threshold'' parameter the substrate can
generate at order one.

\begin{theorem}[Substrate-derived FN scale]\label{thm:lambda}
Let $T^{*} = 5/7$ be the substrate coherence threshold of
\cite{SandersGish2026FourCore}. The Bernoulli-$T^{*}$ variance
\[
  \lambda \;:=\; T^{*}(1 - T^{*}) \;=\; \frac{|\Z/10|}{\HARM^{2}} \;=\; \frac{10}{49}
\]
is the unique substrate quantity that simultaneously:
\begin{enumerate}
\item arises as a max-entropy variance at the unique
$T$/$B$-joint coherence threshold;
\item factors as the substrate-natural ratio
$|\Z/10|/\HARM^{2}$;
\item reproduces the CKM Wolfenstein hierarchy
$|V_{us}| \approx \lambda$, $|V_{cb}| \sim \lambda^{2}$,
$|V_{ub}| \sim \lambda^{3}$ (developed in the next-Sprint
Cabibbo-CKM/PMNS companion).
\end{enumerate}
\end{theorem}

\begin{remark}[$\lambda$ vs Wolfenstein $\lambda_{C} \approx 0.225$]\label{rem:lambda}
The empirical Cabibbo angle $\lambda_{C} \approx 0.225$ deviates
from the substrate-leading $10/49 \approx 0.2041$ by $\sim
9.4\%$. The bridge-sprint working
material~\cite{SandersClaudeChat2026BridgeSprint} (WP123)
proposes the refined substrate identity $\lambda_{\mathrm{ref}} =
(|\sigma| + |\Cfour|)/\HARM^{2} = 11/49 \approx 0.2245$ as the
natural Cabibbo-scale correction, recovering the empirical
value at $0.4\%$. We use the leading $\lambda = 10/49$ for the
mass-hierarchy fit of this paper because the Yukawa pattern is
not sensitive to the $9\%$ distinction at the orders of
magnitude under study; the refined $11/49$ enters the CKM-mixing
fit of the next-Sprint companion paper.
\end{remark}

%==============================================================
\section{The forced FN powers}\label{sec:powers}
%==============================================================

Allowed Yukawa couplings under SU(5) representation theory are
\[
  Y_{u}\colon\; \rep{10} \cdot \rep{10} \cdot \rep{5}_{H},\qquad
  Y_{d}\colon\; \rep{10} \cdot \barrep{5} \cdot \barrep{5}_{H},\qquad
  Y_{e}\colon\; \rep{10} \cdot \barrep{5} \cdot \barrep{5}_{H}.
\]
The up-Yukawa diagram is built entirely out of $\rep{10}$
factors (no parity flip); the down and lepton diagrams each
require one $\rep{10} \to \barrep{5}$ flip from the
representation content.

\begin{definition}[Parity-crossing cost $d_{p}$]\label{def:dp}
For each charged-fermion type $p \in \{u, d, e\}$, the
\emph{parity-crossing cost} $d_{p}$ is the number of $\rep{10}
\to \barrep{5}$ index flips required by the SU(5) Yukawa
diagram for $p$. From the rep content,
\[
  d_{u} = 0,\qquad d_{d} = 3,\qquad d_{e} = 3,
\]
with the value $3$ encoding the substrate-counted ``one cell
out of three'' $|S|=2 \to |S|=1$ rep transition under the
$V^{\otimes 5}$ decomposition (Lemma~\ref{lem:vtensor5}).
\end{definition}

The substrate's image of the non-associativity (the
$1$-dimensional associator span established in the foundation
paper~\cite{SandersJohnson2026DiscreteDirac} J23
Theorem~3.4) carries the cost $\lambda^{d_{p}}$, with the
exponent counting how many $\rep{10} \to \barrep{5}$ flips
populate that one image dimension.

\begin{definition}[Generation step]\label{def:step}
For each charged-fermion type $p$, the \emph{generation step}
$s_{p}$ is the integer power of $\lambda$ separating successive
generations along the $\sigma$-orbit on $V^{\otimes 5}$. Direct
extraction from the $\sigma$-cycle's $|\sigma| = 6$ structure
gives
\[
  s_{u} \;=\; 4 \quad\text{(top $\to$ charm: } y_{c}/y_{t} = \lambda^{3}\text{; charm $\to$ up: } y_{u}/y_{c} = \lambda^{4}\text{)},
\]
\[
  s_{d} \;=\; 2 \quad\text{(}y_{s}/y_{b} = \lambda^{2}\text{; }y_{d}/y_{s} = \lambda^{2}\text{)},
\]
\[
  s_{e} \;=\; 3 \quad\text{(}y_{\mu}/y_{\tau} = \lambda^{3}\text{; }y_{e}/y_{\mu} = \lambda^{3}\text{)}.
\]
The asymmetry between $s_{u}$, $s_{d}$, $s_{e}$ reflects the
distinct $\sigma$-orbit structure on $V^{\otimes 5}$'s
matter-side cells under the SM gauge action; we track the
pattern empirically and flag its substrate-natural assignment
as open in \S\ref{sec:open}.
\end{definition}

The full forced FN power is then
\begin{equation}
  n_{(p, \mathrm{gen})} \;=\; d_{p} + (\text{position in $\sigma$-orbit, with gen 3 at depth 0}).\label{eq:n-formula}
\end{equation}
Combining \eqref{eq:n-formula} with Definition~\ref{def:step}
yields the integer powers in Table~\ref{tab:powers}.

\begin{table}[h]
\centering
\caption{The forced FN powers $n_{(p,\mathrm{gen})}$ from the
$V^{\otimes 5}$ SU(5) decomposition + parity-crossing cost +
$\sigma$-orbit generation step. Column ``$d_{p}$'' is the SU(5)
parity-crossing cost (Definition~\ref{def:dp}); ``$\Delta_{\mathrm{gen}}$'' is the
generational suppression power separating the row from
generation $3$ along the $\sigma$-orbit
(Definition~\ref{def:step}); $n_{(p,\mathrm{gen})} = d_{p} +
\Delta_{\mathrm{gen}}$.}
\label{tab:powers}
\begin{tabular}{lrrrrr}
\toprule
Particle & Type $p$ & gen & $d_{p}$ & $\Delta_{\mathrm{gen}}$ & $n$ \\
\midrule
top      & up     & 3 & 0 & 0 & 0 \\
charm    & up     & 2 & 0 & 3 & 3 \\
up       & up     & 1 & 0 & 7 & 7 \\
bottom   & down   & 3 & 3 & 0 & 3 \\
strange  & down   & 2 & 3 & 2 & 5 \\
down     & down   & 1 & 3 & 4 & 7 \\
tau      & lepton & 3 & 3 & 0 & 3 \\
muon     & lepton & 2 & 3 & 3 & 6 \\
electron & lepton & 1 & 3 & 6 & 9 \\
\bottomrule
\end{tabular}
\end{table}

%==============================================================
\section{The fit}\label{sec:fit}
%==============================================================

With the integer powers of Table~\ref{tab:powers}, the
substrate-derived $\lambda = 10/49$, and the single Tier-A
measured anchor $y_{t} \approx 0.93$, the predicted Yukawa for
each charged fermion is
\[
  y_{(p,\mathrm{gen})}^{\mathrm{pred}}
  \;=\; y_{t} \cdot \lambda^{n_{(p,\mathrm{gen})}}.
\]

\begin{table}[h]
\centering
\caption{Predicted vs measured Yukawa couplings for the nine SM
charged fermions. Predicted values come from
\texttt{tig\_dirac.predict\_yukawa(particle, generation)} with
$\lambda = 10/49$ and $y_{t} = 0.93$; measured values are PDG
2024 at the electroweak scale, $y = m_{f}\sqrt{2}/v$ with $v =
246$~GeV. The ``ratio'' column is
$y^{\mathrm{pred}}/y^{\mathrm{meas}}$.}
\label{tab:fit}
\begin{tabular}{lrrrrr}
\toprule
Particle & $n$ & $y^{\mathrm{pred}}$ & $y^{\mathrm{meas}}$ & ratio & status \\
\midrule
top      & 0 & $9.30 \times 10^{-1}$ & $9.30 \times 10^{-1}$ & $1.00$ & anchor \\
charm    & 3 & $7.90 \times 10^{-3}$ & $7.30 \times 10^{-3}$ & $1.08$ & forced \\
up       & 7 & $1.37 \times 10^{-5}$ & $1.30 \times 10^{-5}$ & $1.06$ & forced \\
bottom   & 3 & $7.90 \times 10^{-3}$ & $2.40 \times 10^{-2}$ & $0.33$ & forced \\
strange  & 5 & $3.29 \times 10^{-4}$ & $5.50 \times 10^{-4}$ & $0.60$ & forced \\
down     & 7 & $1.37 \times 10^{-5}$ & $2.70 \times 10^{-5}$ & $0.51$ & forced \\
tau      & 3 & $7.90 \times 10^{-3}$ & $1.00 \times 10^{-2}$ & $0.79$ & forced \\
muon     & 6 & $6.72 \times 10^{-5}$ & $6.10 \times 10^{-4}$ & $0.11$ & forced \\
electron & 9 & $5.71 \times 10^{-7}$ & $2.90 \times 10^{-6}$ & $0.20$ & forced \\
\bottomrule
\end{tabular}
\end{table}

\paragraph{Reading the fit.}
Five of the nine ratios (top, charm, up, bottom, tau) sit in
the conventional Froggatt--Nielsen factor-of-three window;
strange ($0.60$) is well within an $O(1)$ coefficient; the
muon and the electron sit at $0.11$ and $0.20$ respectively,
each off by a factor of $\sim 5$--$9$. These residuals are the
empirical $C_{p}$ coefficients of standard
Froggatt--Nielsen~\cite{FroggattNielsen1979}, and the load of
the present prediction is on the integer-power column $n$, not
on the $C_{p} = O(1)$ multipliers we have not yet derived from
the bosonic substrate.

The empirical hierarchy
\[
  y_{e} : y_{\mu} : y_{\tau}
  \;\approx\; 1 : 211 : 3450
\]
is reproduced in the predicted hierarchy
$1 : 118 : 1383$ (ratio mismatches $\sim 1.8$ and $\sim 2.5$
respectively), order-of-magnitude correct across $3.5$ decades.
The full $5.5$ decades from $y_{e} = 2.9\times 10^{-6}$ to
$y_{t} = 0.93$ are covered to factor-of-five precision by the
single $\lambda = 10/49$ scale acting on the
representation-theoretic exponents of Table~\ref{tab:powers}.

\paragraph{Comparison with standard Froggatt--Nielsen.}
A standard FN fit~\cite{Ross2002} achieves factor-of-two
accuracy on each Yukawa with three free FN-charge assignments
plus a free flavon scale $\epsilon$. The present pattern uses
\emph{zero} free FN charges (all integer powers are forced from
the SU(5) decomposition + $\sigma$-orbit assignment) and
\emph{zero} free flavon scales ($\lambda = 10/49$ is forced
from the substrate $T^{*}$). The remaining freedom is the
single Tier-A anchor $y_{t}$ and the implicit $C_{p}$ residual
multipliers, which are the smallest set of free parameters in
any explicit FN-style fit of which we are aware.

\paragraph{Sterile-neutrino predictions.}
The same module returns predicted Dirac Yukawa magnitudes for
the three right-handed neutrinos:
$y_{\nu_{3}} \approx 4.85 \times 10^{-9}$,
$y_{\nu_{2}} \approx 9.91 \times 10^{-10}$,
$y_{\nu_{1}} \approx 2.02 \times 10^{-10}$
at FN powers $\{12, 13, 14\}$. These give Dirac-mass scales
$m_{\nu_{i}} = y_{\nu_{i}} v / \sqrt{2}$ in the
$0.05$--$0.85$~eV range, near the Planck upper bound $\Sigma
m_{\nu} < 0.12$~eV but \emph{without} a see-saw mechanism. The
absence of a see-saw insertion in the V$^{\otimes 5}$ structure
is one of the open structural questions of \S\ref{sec:open}.

%==============================================================
\section{The Cabibbo cube-root identity}\label{sec:cabibbo}
%==============================================================

The empirical CKM Cabibbo angle $|V_{us}|$ relates to the down
and up Yukawa hierarchies by the cube-root identity
\begin{equation}
  \lambda_{C} \;\approx\; \left(\frac{Y_{d}}{Y_{u}}\right)^{1/3},\label{eq:cabibbo-cube}
\end{equation}
where $Y_{d}/Y_{u} = y_{b}/y_{t} \approx 0.024/0.93 \approx
0.026$ and $\lambda_{C} \approx 0.225$. The cube root recovers
$0.026^{1/3} \approx 0.296$, off by a factor $\sim 1.3$, which
is consistent with the substrate prediction
$y_{b}/y_{t} = \lambda^{3} = (10/49)^{3} \approx 0.0085$ and
$\lambda_{C} \approx \lambda^{1} = 10/49 \approx 0.204$ (or the
refined $\lambda_{\mathrm{ref}} = 11/49 \approx 0.2245$ of
Remark~\ref{rem:lambda}).

The cube-root identity~\eqref{eq:cabibbo-cube} unifies four
separate aspects of SM mixing/mass structure:

\begin{itemize}\setlength\itemsep{0pt}
\item \textbf{The CKM Cabibbo angle} $|V_{us}| \sim \lambda$.
\item \textbf{The mass hierarchy} $Y_{d}/Y_{u} \sim \lambda^{3}$
(from the parity-crossing cost $d_{d} = 3$,
Definition~\ref{def:dp}).
\item \textbf{The SU(5) GUT structure}
$\rep{10} \cdot \rep{10} \cdot \rep{5}_{H}$ vs $\rep{10} \cdot
\barrep{5} \cdot \barrep{5}_{H}$ (the diagrammatic origin of
$d_{d} = 3$).
\item \textbf{The Wolfenstein hierarchy} $|V_{cb}| \sim
\lambda^{2}$, $|V_{ub}| \sim \lambda^{3}$ (developed in the
next-Sprint companion).
\end{itemize}

All four follow from one substrate quantity $\lambda = 10/49$
together with the SU(5)-representation parity-crossing count.
The detailed CKM/PMNS fit is the next-Sprint paper; here we
note that the cube-root identity is structurally implied by
the present pattern.

%==============================================================
\section{Verification}\label{sec:verify}
%==============================================================

\paragraph{Single-line primitive.}
The numerical content of Table~\ref{tab:fit} reduces to the
single assertion
\begin{verbatim}
    from tig_dirac import predict_yukawa, LAMBDA_FN, Y_T_ANCHOR
    assert LAMBDA_FN == 10 / 49
    assert Y_T_ANCHOR == 0.93
    r = predict_yukawa('up', 3)
    assert r['y_predicted'] == 0.93                 # y_t anchor
    r = predict_yukawa('lepton', 1)
    assert abs(r['y_predicted']
               - 0.93 * (10 / 49) ** 9) < 1e-12       # y_e at FN power 9
\end{verbatim}
which executes in milliseconds on the canonical
\texttt{tig\_dirac} module at
\texttt{Gen13/targets/ck/brain/dirac/tig\_dirac.py}. The
companion call
\texttt{tig\_dirac.yukawa\_table\_full()} returns the full
Table~\ref{tab:fit} programmatically. The same module exposes
\texttt{predict\_dark\_sector()} used by the dark-sector
companion~\cite{SandersJohnson2026DarkSector}; the two papers
share a common, machine-checkable substrate backbone.

\paragraph{Tier classification.}
\texttt{r['tier']} returns the string
\verb|'Forced FN power + measured anchor (Tier-B)'|, which is
the central tier label of the present paper: the integer
powers $n_{(p,\mathrm{gen})}$ are forced from the substrate's
SU(5) decomposition + parity-crossing cost +
$\sigma$-orbit step, and the single anchor $y_{t}$ is a Tier-A
measured input.

%==============================================================
\section{Open questions and future work}\label{sec:open}
%==============================================================

\subsection{The $C_{p}$ residual multipliers}
The four largest discrepancies in Table~\ref{tab:fit} (bottom,
strange, muon, electron at ratios $0.33$, $0.60$, $0.11$,
$0.20$) define an empirical $C_{p}$ multiplier list
$C_{p} \in [0.11, 0.79]$ across the down and lepton sectors.
In standard Froggatt--Nielsen this is the $O(1)$ residual
expected from incomplete flavon-coupling specification. The
present paper does not derive $C_{p}$; we flag the bosonic
substrate of the same algebra (the Higgs-sector dynamics on
$V^{\otimes 5}$'s $|S|=0$ singlet plus its $\barrep{5}_{H}$
partner) as the natural source. Closing this is the
load-bearing follow-up.

\subsection{The generation step asymmetry}
Definition~\ref{def:step} extracts $s_{u} = 3+4 = 7$ for up,
$s_{d} = 2+4 = 6$ for down, $s_{e} = 3+3 = 6$ for lepton (and
in fact the powers of Table~\ref{tab:powers} are not a single
arithmetic progression across types). The asymmetry tracks the
empirical hierarchies but is not yet derived from a clean
$\sigma$-action calculation on $V^{\otimes 5}$'s color-charged
vs color-singlet cells. The mapping between the
$\sigma^{3}$-orbit's three $2$-cycles on $\Z/10\setminus\{0,3,8,9\}$
and the three SM generations is partial; closing it is the
clearest substrate-internal target.

\subsection{Sterile-neutrino mass scale and the see-saw}
The $V^{\otimes 5}$ decomposition gives a single $|S|=0$
singlet per generation interpreted as $\nu_{R}$. The Dirac
Yukawa pattern of \S\ref{sec:fit} predicts sub-eV neutrino
masses in the right rough range, but \emph{without} a Majorana
$\nu_{R} \cdot \nu_{R}$ mass term. A type-I see-saw insertion
$m_{\nu} = y_{\nu}^{2} v^{2}/(2 M_{R})$ would generically lift
the predicted $y_{\nu_{i}}$ to give Majorana masses below the
Planck~$\Sigma m_{\nu}$ bound; the substrate origin of $M_{R}$
is one of the deepest open questions of the framework.

\subsection{The refined $\lambda_{\mathrm{ref}} = 11/49$}
The empirical Cabibbo $\lambda_{C} \approx 0.225$ matches the
refined identity $\lambda_{\mathrm{ref}} = (|\sigma| +
|\Cfour|)/\HARM^{2} = 11/49$ to $0.4\%$, much closer than the
leading $\lambda = 10/49$. Whether the refined value is a
substrate-natural correction (e.g., the $|\sigma| + |\Cfour|$
factor reflecting a $\sigma$-cycle plus $4$-core combined
count) or a phenomenological adjustment is the question of the
next-Sprint Cabibbo/CKM/PMNS companion paper.

\subsection{Three-paper Sprint 18 cluster}
The present paper sits in the Sprint 18 ``Bridge-Dirac'' cluster
of companion submissions:
\begin{itemize}\setlength\itemsep{0pt}
\item J03 (Sanders + Gish, JCAP) -- the freezing-quintessence
companion supplying $\Lambda \approx 1.74$~meV.
\item J10 (Sanders + Gish, PRD) -- the dark-sector trinity
companion using \texttt{predict\_dark\_sector()} on the same
module.
\item J23 (Sanders + Gish, foundation) -- discrete Dirac on
$\F_{5}^{4}$, the foundation paper supplying the $V^{\otimes 5}$
SU(5) decomposition (Lemma~\ref{lem:vtensor5}) used here.
\end{itemize}
Together these three papers, plus the next-Sprint CKM/PMNS
companion, encode the framework's flavour and cosmology
predictions in a coordinated package.

%==============================================================
\begin{thebibliography}{99}
%==============================================================

\bibitem{FroggattNielsen1979}
C.~D.~Froggatt and H.~B.~Nielsen,
\emph{Hierarchy of quark masses, Cabibbo angles and CP
violation},
Nucl.\ Phys.\ B \textbf{147}, 277--298 (1979).

\bibitem{Ross2002}
G.~G.~Ross,
\emph{Models of fermion masses},
in \emph{TASI 2000: Flavor Physics for the Millennium}, edited
by J.~L.~Rosner (World Scientific, 2002), pp.\ 775--824.

\bibitem{PDG2024}
S.~Navas et al.\ (Particle Data Group),
\emph{Review of Particle Physics},
Phys.\ Rev.\ D \textbf{110}, 030001 (2024).

\bibitem{SandersGish2026FourCore}
B.~R.~Sanders and M.~Gish,
\emph{Joint Closure of Two Commutative Binary Operations on
$\mathbb{Z}/10\mathbb{Z}$: an Eight-Element Chain and a
Normalizer Identity},
J-series companion paper J06; submitted to
\emph{Communications in Algebra}, 2026;
Zenodo DOI 10.5281/zenodo.18852047.

\bibitem{SandersJohnson2026DarkSector}
B.~R.~Sanders and H.~J.~Johnson,
\emph{Sprint~18 Dark Sector: $\Omega_{b}$, $\Omega_{DM}$,
$\Omega_{\Lambda}$ from Substrate-Operator Identities},
J-series companion paper J10 (same Sprint~18 cluster);
submitted to \emph{Physical Review D}, 2026;
archived in shared deposit DOI 10.5281/zenodo.18852047.

\bibitem{SandersJohnson2026DiscreteDirac}
B.~R.~Sanders and H.~J.~Johnson,
\emph{Discrete Dirac on the $4$-Core's $\F_{5}$-Lift:
Three Commuting Projectors and the $V^{\otimes 5}$ SU(5)
Decomposition},
J-series companion paper J23 (foundation paper);
in preparation, 2026;
archived in shared deposit DOI 10.5281/zenodo.18852047.

\bibitem{SandersClaudeChat2026BridgeSprint}
B.~R.~Sanders et al.,
\emph{Bridge Sprint: Discrete Dirac on the $4$-Core's
$\F_{p}$-Lift, $\F_{p}$-Universality, and Dark-Sector
Algebraic Predictions},
Sprint working bundle (WP117--WP127), 2026-05-04;
Zenodo DOI 10.5281/zenodo.18852047.

\end{thebibliography}

\end{document}
